\documentclass[11pt]{article}
\usepackage[a4paper,margin=1in]{geometry}
\usepackage{fontspec}
\usepackage[boldfont=false]{xeCJK}
\setCJKmainfont{FandolSong-Regular.otf}
\usepackage{amsmath}
\usepackage{amssymb}
\usepackage{array}
\usepackage{booktabs}
\usepackage{graphicx}
\usepackage{adjustbox}
\usepackage{enumitem}
\setlist[itemize]{topsep=2pt,partopsep=0pt,parsep=0pt,itemsep=2pt,leftmargin=1.4em}
\setlist[enumerate]{topsep=2pt,partopsep=0pt,parsep=0pt,itemsep=2pt,leftmargin=1.6em}
\usepackage{parskip}
\usepackage{fvextra}
\fvset{breaklines=true,breakanywhere=true,fontsize=\small}
\setlength{\parindent}{0pt}

\begin{document}

\begin{center}
  {\LARGE\bfseries Chemistry of transition elements}\\[0.35em]
  {\large A Level Chemistry}
\end{center}
\vspace{0.6em}

\section*{Transition elements}

A \textbf{transition element} 过渡元素 is a d-block element that forms one or more stable ions with \textbf{incomplete d orbitals}. (Scandium and zinc are in the d-block but are not transition elements, because their stable ions have empty or full d orbitals.)

The 3d \textbf{orbitals} 轨道 have set shapes: the $3\text{d}_{xy}$ orbital has four lobes pointing between the axes, and the $3\text{d}_{z^2}$ orbital has two lobes along the $z$-axis with a ring around the middle.

\subsection*{Four key properties (and why)}

\begin{center}
\adjustbox{max width=\linewidth}{%
\begin{tabular}{ll}
\toprule
\textbf{Property} & \textbf{Reason} \\
\midrule
variable \textbf{oxidation state} 氧化态 & the 3d and 4s sub-shells are close in energy, so similar small amounts of energy remove different numbers of electrons \\
act as a \textbf{catalyst} 催化剂 & they have more than one stable oxidation state, and vacant d orbitals that can form dative bonds \\
form complex ions & vacant d orbitals can accept lone pairs \\
form coloured compounds & electrons move between split d orbitals (see below) \\
\bottomrule
\end{tabular}}
\end{center}

\section*{Ligands and complexes}

A transition metal ion can be surrounded by \textbf{complex ions} 配离子. The species attached are ligands.

A \textbf{ligand} 配体 is a species with a lone pair of electrons that forms a \textbf{dative covalent bond} 配位键 to the central metal ion. (The \textbf{lone pair} 孤对电子 is what it donates.) Ligands are grouped by how many such bonds they can form:

\begin{itemize}
\item \textbf{monodentate} 单齿: one bond ($\text{H}_2\text{O}$, $\text{NH}_3$, $\text{Cl}^-$, $\text{CN}^-$).

\item \textbf{bidentate} 双齿: two bonds (1,2-diaminoethane "en", and the ethanedioate ion $\text{C}_2\text{O}_4^{2-}$).

\item \textbf{polydentate} 多齿: many bonds ($\text{EDTA}^{4-}$, which uses six).

\end{itemize}

A \textbf{complex} 配合物 is a central metal atom or ion surrounded by one or more ligands. Its shape can be \textbf{linear} 直线形, \textbf{square planar} 平面正方形, \textbf{tetrahedral} 四面体形 or \textbf{octahedral} 八面体形.

The \textbf{coordination number} 配位数 is the number of dative bonds from the ligands to the central ion (6 → octahedral, 4 → tetrahedral or square planar, 2 → linear). To predict the charge of a complex, add the metal's charge and all the ligand charges.

\subsection*{Ligand exchange}

In \textbf{ligand exchange} 配体交换 one ligand replaces another, often with a colour change. For copper(II):

\[
[\text{Cu}(\text{H}_2\text{O})_6]^{2+} \;(\text{pale blue}) \;\xrightarrow{\text{NH}_3}\; [\text{Cu}(\text{NH}_3)_4(\text{H}_2\text{O})_2]^{2+} \;(\text{deep blue})
\]

With concentrated $\text{HCl}$ it becomes yellow $[\text{CuCl}_4]^{2-}$; cobalt(II) behaves in a similar way.

\subsection*{Redox reactions of transition ions}

Use $E^{\ominus}$ values to predict whether a redox reaction is feasible. Common titrations you should be able to calculate include $\text{MnO}_4^-/\text{C}_2\text{O}_4^{2-}$ and $\text{MnO}_4^-/\text{Fe}^{2+}$ in acid (purple to colourless), and $\text{Cu}^{2+}/\text{I}^-$ (which makes iodine, then titrated with thiosulfate).

\section*{Why complexes are coloured}

In a free ion the five d orbitals are \textbf{degenerate} 简并 — they have the same energy. When ligands come close, they split the d orbitals into two \textbf{non-degenerate} 非简并 sets, separated by an energy gap $\Delta E$:

\begin{itemize}
\item \textbf{octahedral}: three lower and two higher orbitals.

\item \textbf{tetrahedral}: two lower and three higher orbitals.

\end{itemize}

A complex absorbs light whose frequency matches $\Delta E$, promoting an electron from a lower to a higher d orbital. The colour you \textbf{see} is the \textbf{complementary colour} 互补色 of the light absorbed. Different ligands give a different $\Delta E$, so they change the frequency absorbed and hence the colour — which is why ligand exchange changes the colour.

\section*{Stereoisomerism in complexes}

\begin{itemize}
\item \textbf{geometrical isomerism} 几何异构 (\textbf{cis} 顺式 / \textbf{trans} 反式) appears in square planar complexes such as $[\text{Pt}(\text{NH}_3)_2\text{Cl}_2]$, and in octahedral complexes such as $[\text{Co}(\text{NH}_3)_4(\text{H}_2\text{O})_2]^{2+}$.

\item \textbf{optical isomerism} 旋光异构 appears in octahedral complexes with bidentate ligands, such as $[\text{Ni}(\text{en})_3]^{2+}$, which has two non-superimposable mirror images.

\end{itemize}

You can also deduce the polarity of a complex: a \emph{cis} form may be polar, while the matching \emph{trans} form is often non-polar because its dipoles cancel.

\section*{Stability constants}

The \textbf{stability constant} 稳定常数 ($K_{\text{stab}}$) is the equilibrium constant for forming a complex ion from the metal ion and its ligands in solution (water is left out of the expression).

A \textbf{large} $K_{\text{stab}}$ means a very stable complex. In a ligand exchange, the position moves towards the complex with the larger $K_{\text{stab}}$ — that is why a ligand that forms a more stable complex can push out a weaker one.


\end{document}
